\section{Instruction sets}
\label{sec:isa}

Pollard rho on these curves needs two kinds of arithmetic: products of
binary polynomials (for $\FF_{2^{131}}$) and multi-word integer products
with long carry chains (for the NIST primes). Both have gained dedicated
instructions over the last fifteen years, on GPUs and CPUs alike.
Table~\ref{tab:isa} and Figure~\ref{fig:isa} summarise the ones this
paper uses or prices; the rest of this section says why each matters.

\begin{table}[t]
\centering
\footnotesize
\setlength{\tabcolsep}{4pt}
\caption{Instructions relevant to rho on $\FF_{2^{131}}$ and on prime
  fields, with the first generation that has each. ``Use here'' says how
  this paper's code uses it: in the measured clients, priced, or not
  yet tried.}
\label{tab:isa}
\begin{tabular}{@{}p{0.2\textwidth}p{0.31\textwidth}p{0.2\textwidth}p{0.2\textwidth}@{}}
\toprule
instruction & what it does & first available & use here \\
\midrule
\multicolumn{4}{@{}l}{\emph{NVIDIA GPUs}} \\
\texttt{PRMT}, \texttt{POPC}, \texttt{FLO} & byte permute; population count; find leading one & Fermi (2010) & every client \\
\texttt{SHF} & funnel shift of a register pair & Kepler (2012--13) & every client \\
\texttt{LOP3} & any 3-input Boolean function & Maxwell (2014) & every client \\
\texttt{IADD3}, native \texttt{IMAD} & 3-input add; 32-bit multiply-add & Maxwell; Volta (2017) & NIST walk \\
\texttt{mad.lo.cc}, \texttt{madc.hi.cc} & multiply-add with carry in/out (PTX) & all CUDA GPUs & NIST walk \\
\texttt{clmad.lo}/\texttt{.hi} & $64{\times}64$ carry-less multiply-add (PTX 9.3) & CUDA 13.3 (2026), \texttt{sm\_80}+ & ECC2K-130 client \\
\midrule
\multicolumn{4}{@{}l}{\emph{x86 (Intel, AMD)}} \\
\texttt{PCLMULQDQ} & $64{\times}64$ carry-less multiply & Westmere (2010) & packed CPU client \\
\texttt{MULX}, \texttt{ADCX}/\texttt{ADOX} & flag-free multiply; two carry chains & Haswell (2013); Broadwell (2014) & not tried \\
\texttt{VPTERNLOGD} & any 3-input Boolean function, 512-bit & Skylake-SP (2017) & bitsliced client \\
\texttt{VPMADD52} (IFMA) & $52{\times}52$-bit multiply-add & Cannon Lake (2018) & not tried \\
\texttt{VPCLMULQDQ}, GFNI & 4 carry-less products per 512 bits; GF($2^8$) affine & Ice Lake (2019); AMD Zen 4 (2022) & priced \\
\midrule
\multicolumn{4}{@{}l}{\emph{Arm}} \\
\texttt{PMULL}/\texttt{PMULL2} & $64{\times}64$ carry-less multiply & Armv8.0 Crypto & packed CPU client \\
\texttt{EOR3}, \texttt{BCAX} & 3-way XOR; AND-NOT then XOR & Armv8.2 SHA3 & bitsliced client \\
SVE2 \texttt{PMULLB/T}, \texttt{EOR3} & vector carry-less and 3-input logic & Armv9.0 (SVE2-AES) & not tried \\
\midrule
\multicolumn{4}{@{}l}{\emph{RISC-V}} \\
Zbc, Zvbc & scalar and vector carry-less multiply & ratified 2021, 2023 & not tried \\
\bottomrule
\end{tabular}
\end{table}

\paragraph{NVIDIA GPUs.} The integer pipe has grown three instructions
that a binary field uses on every step. \texttt{PRMT} and \texttt{POPC}
date from Fermi, the funnel shift \texttt{SHF} from Kepler, and
\texttt{LOP3}, which evaluates any Boolean function of three inputs from an
8-bit truth table, from Maxwell~\cite{nvidia_binutils}. \texttt{LOP3}
alone is worth $1.65\times$ on the bitsliced product
(Section~\ref{sec:arith-bitsliced}). \texttt{POPC} still issues at a
quarter of the logic rate on every current part, 16 per clock per SM, so
each Hamming weight costs about four logic slots~\cite{nvidia_bpg134}.
The prime-field side relies on PTX's carry-chain forms, which every CUDA
GPU has; the extended-precision forms issue at half the rate of a plain
multiply-add~\cite{nvidia_bpg134}.

The change that matters most for ECC2K-130 is recent. PTX ISA~9.3,
shipped with CUDA~13.3 in May 2026, added \texttt{clmad.lo} and
\texttt{clmad.hi}: the low and high halves of a $64\times64$ carry-less
product, XORed into an addend, for \texttt{sm\_80} and
later~\cite{nvidia_ptx93}. NVIDIA describes it as accelerated in hardware
since Ampere and reports GHASH at about 6.3~TB/s on a B200 and 1.3~TB/s
on an RTX~5090~\cite{watson2026clmad}. NVIDIA documents neither its SASS
mnemonic nor its issue rate, which is why we measured both
(Section~\ref{sec:blackwell-rates}). Nsight attributes it to the FP64
pipe, and it is slow exactly where FP64 is slow: about two lanes per
SM-clock on the RTX~PRO~6000, ten to eighteen times that on the B200 and
H100. That matters because
NVIDIA's FP64 ratio varies by product rather than by generation: $2{:}1$
against FP32 on the A100, H100 and B200, but $64{:}1$ on the GA10x and Ada
consumer parts, on the RTX~PRO~6000 and on the
B300~\cite{nvidia_ga102_wp,nvidia_ada_wp,nvidia_h100_wp,nvidia_rtx_blackwell_wp,aubrey2025blackwellultra}.

Two other generational changes show up in our rates. From Ampere's GA10x
parts onward each SM partition pairs 16 FP32-only lanes with 16 lanes that
do FP32 or INT32, and Ada kept that split; the workstation Blackwell
chips (\texttt{sm\_120}) unify them, so every lane can issue integer
work~\cite{nvidia_ga102_wp,nvidia_ada_wp,nvidia_rtx_blackwell_wp}. That is
consistent with the factor of two we measure per SM between Ada and
\texttt{sm\_120} on this integer-bound walk (Section~\ref{sec:blackwell-logic}).
Hopper's DPX instructions (fused add with min or max for dynamic
programming), its tensor memory accelerator and Blackwell's
\texttt{tcgen05} tensor-memory instructions do not help this walk; the
products here have no operand shared across lanes, which tensor cores
need (Section~\ref{sec:open-method}).

\paragraph{x86.} Intel added \texttt{PCLMULQDQ} with Westmere in
2010~\cite{gueron2014clmul}, \texttt{MULX} with Haswell and the dual carry
chains \texttt{ADCX}/\texttt{ADOX} with Broadwell~\cite{ozturk2012adx},
AVX-512 with its ternary-logic \texttt{VPTERNLOGD} in Skylake-SP, the
52-bit integer multiply-add IFMA in Cannon Lake~\cite{gueron2016ifma}, and
the 512-bit \texttt{VPCLMULQDQ} (four carry-less products per instruction)
with GFNI in Ice Lake~\cite{intel_sdm_vol1,drucker2018vpclmulqdq}. AMD
implemented AVX-512, including \texttt{VPCLMULQDQ} and GFNI, in Zen~4 on a
256-bit datapath that issues each 512-bit operation in two halves, and
Zen~5's server parts run it at full width~\cite{amd_58003,amd_epyc5_wp}.
Intel's AVX10.2 makes 512-bit vectors universal and APX doubles the
general-purpose registers to 32~\cite{intel_avx102,intel_apx}; both would
help a register-starved bitsliced leaf (Section~\ref{sec:cpu-bitsliced}).

\paragraph{Arm and RISC-V.} Armv8.0's optional Crypto extension has the
$64\times64$ carry-less \texttt{PMULL}; the Armv8.2 SHA3 extension adds the
three-input \texttt{EOR3} and \texttt{BCAX}; SVE2 brings vector forms of
both, with 128-bit-result \texttt{PMULLB}/\texttt{PMULLT} in the SVE2-AES
extension~\cite{arm_a64_2026}. Graviton2 (Neoverse N1) has PMULL but not
SHA3, Graviton3 (V1) adds SHA3 and SVE, and Graviton4 (V2) has SVE2 with
the 128-bit PMULL; Apple's M4 has SHA3 and SME but no non-streaming SVE.
RISC-V ratified scalar carry-less multiply (Zbc) in 2021 and its vector
form (Zvbc) in 2023~\cite{riscv_bitmanip,riscv_vcrypto}.

\begin{figure}[t]
\centering
\begin{tikzpicture}[x=0.78cm,y=1cm,font=\scriptsize,
  ev/.style={circle,inner sep=1.7pt},
  lbl/.style={align=center,font=\scriptsize}]
% axis
\draw[inkS,thick] (0,0) -- (16.6,0);
\foreach \y in {0,2,...,16}{
  \draw[inkS] (\y,-0.08) -- (\y,0.08);
  \node[below,font=\scriptsize] at (\y,-0.1) {\the\numexpr2010+\y\relax};
}
\node[anchor=east,font=\footnotesize\bfseries] at (-0.2,1.6) {NVIDIA};
\node[anchor=east,font=\footnotesize\bfseries] at (-0.2,-1.9) {CPUs};
% GPU events (above)
\foreach \x/\h/\t/\c in {
  0/0.7/{\texttt{PRMT}, \texttt{POPC}\\(Fermi)}/cA,
  2.6/1.55/{\texttt{SHF}\\(Kepler)}/cA,
  4.2/0.7/{\texttt{LOP3}, \texttt{IADD3}\\(Maxwell)}/cA,
  7.4/1.55/{native \texttt{IMAD}\\(Volta)}/cA,
  10.4/0.7/{INT32/FP32 split\\(GA10x)}/cA,
  12.7/1.55/{DPX, clusters\\(Hopper)}/cA,
  15/0.7/{unified INT32\\(\texttt{sm\_120})}/cA,
  16.4/1.55/{\texttt{clmad}\\(CUDA 13.3)}/cB}{
  \draw[gridc] (\x,0) -- (\x,\h-0.2);
  \node[ev,fill=\c] at (\x,0) {};
  \node[lbl] at (\x,\h) {\t};
}
% CPU events (below)
\foreach \x/\h/\t/\c in {
  0/-1.0/{\texttt{PCLMULQDQ}\\(Westmere)}/cB,
  1.6/-1.9/{\texttt{PMULL}\\(Armv8)}/cB,
  3.5/-1.0/{\texttt{MULX}\\(Haswell)}/cA,
  5.4/-1.9/{\texttt{ADX}; \texttt{EOR3}\\(Broadwell; Armv8.2)}/cA,
  7.3/-1.0/{\texttt{VPTERNLOGD}\\(Skylake-SP)}/cA,
  9.2/-1.9/{\texttt{VPCLMULQDQ}, GFNI\\(Ice Lake)}/cB,
  11.3/-1.0/{SVE2; Zbc\\(Armv9; RISC-V)}/cB,
  13.0/-1.9/{AVX-512\\(AMD Zen 4)}/cA,
  15.2/-1.0/{full-width AVX-512\\(Zen 5)}/cA}{
  \draw[gridc] (\x,0) -- (\x,\h+0.25);
  \node[ev,fill=\c] at (\x,0) {};
  \node[lbl] at (\x,\h) {\t};
}
\node[ev,fill=cB] at (0.2,-2.75) {}; \node[anchor=west] at (0.4,-2.75) {carry-less multiply};
\node[ev,fill=cA] at (4.6,-2.75) {}; \node[anchor=west] at (4.8,-2.75) {logic, shifts and integer carry chains};
\end{tikzpicture}
\caption{When the instructions this arithmetic relies on arrived. Dates
  are first availability in a shipping product (Armv8 and Armv8.2 are
  architecture versions; Zbc is the ratification date). The GPU
  carry-less multiply arrived last of all, in 2026, which is why the
  ECC2K-130 GPU client is new work rather than a port.}
\label{fig:isa}
\end{figure}

